\documentclass[12pt]{article}

\usepackage[margin=1in]{geometry}
\usepackage{booktabs}
\usepackage{makecell}
\usepackage{setspace}
\onehalfspacing

\title{Information and Support for Government Intervention:\\
Evidence from the 2015 Swiss Franc Shock}
\author{Your Name}
\date{\today}

\begin{document}
\maketitle

\section{Data}

We use a survey of 2{,}044 Polish adults fielded after the January 2015 removal of the Swiss franc exchange-rate floor. Respondents were randomly assigned to one of four conditions before being asked whether the government should intervene to help borrowers with foreign-currency loans: a control group with no information, a factual information group, a historical context group, and a group shown the Hungarian policy response. The outcome, \texttt{supports\_intervention}, equals one if the respondent supports some or big intervention and zero otherwise. Survey weights are used where noted.

\section{Results}

% TODO (exercise step 3): the results paragraph goes here, written by Claude Code from the
% table below. Do not type numbers into this file yourself.

Table~\ref{tab:main} reports logit coefficients for the effect of each information condition relative to the control group.

\begin{table}[htbp]
  \centering
  \caption{Effect of Information Conditions on Support for Intervention}
  \label{tab:main}
  \setlength{\tabcolsep}{4pt}
\begin{tabular}{lcccccc}
\toprule
& (1) Base & (2) Unwtd. & (3) Basic & (4) Full & (5) Any & (6) Urban \\
\midrule
Information & \makecell{0.302$^{*}$\\ (0.126)} & \makecell{0.290$^{*}$\\ (0.126)} & \makecell{0.324$^{*}$\\ (0.127)} & \makecell{0.267\\ (0.170)} &  & \makecell{0.326$^{+}$\\ (0.167)} \\
History & \makecell{0.123\\ (0.128)} & \makecell{0.047\\ (0.127)} & \makecell{0.133\\ (0.130)} & \makecell{0.261\\ (0.174)} &  & \makecell{0.242\\ (0.169)} \\
Hungary comparison & \makecell{0.194\\ (0.130)} & \makecell{0.145\\ (0.129)} & \makecell{0.185\\ (0.131)} & \makecell{0.225\\ (0.176)} &  & \makecell{0.221\\ (0.169)} \\
Any information &  &  &  &  & \makecell{0.218$^{*}$\\ (0.107)} &  \\
\midrule
N & 2,036 & 2,036 & 2,035 & 1,151 & 2,035 & 1,208 \\
Pseudo $R^2$ & 0.002 & 0.002 & 0.019 & 0.023 & 0.018 & 0.026 \\
\bottomrule
\end{tabular}

\vspace{0.5em}

{\footnotesize Logit coefficients (log-odds), model-based standard errors in parentheses. Reference arm: control. Columns (1) and (3)--(6) use survey weights; column (2) does not. Columns (3)--(6) control for FX status, age, female, education, and settlement size; column (4) adds income quintile and left-right, which restricts the sample to respondents with both. Column (5) replaces the three arms with a single any-information indicator. Column (6) keeps respondents in small or big urban areas only. \par}

{\footnotesize $^{***}p<0.001$, $^{**}p<0.01$, $^{*}p<0.05$, $^{+}p<0.1$ \par}

\end{table}

\end{document}
